\documentclass{ximera}

\begin{document}

    \title{In-Class Assignment 5}
    
    \begin{abstract} Solve the following problems.
    \end{abstract}
    \maketitle


Let's consider the system of equations 
\begin{align*}
    x+2y-z &= 3 \\
    2x+5y+z &=12 \\
    3x+8y+2z &=19
\end{align*}
which, to allow us to perform Gaussian elimination, can be written in the form:
\begin{align*}
\left[
\begin{array}{rrr|r}
1 & 2 & -1 & 3 \\
2 & 5 & 1 & 12 \\
3 & 8 & 2 & 19
\end{array}
\right],
\end{align*}
where $A = \begin{bmatrix} 1 & 2 & -1 \\
2 & 5 & 1 \\
3 & 8 & 2\end{bmatrix}$.

\begin{exercise}
    Determine the row operations necessary to put $A$ into an upper triangular form.  Write these row operations as elementary matrices multiplying $A$ (there should be 3).  
\end{exercise}
\vfill
\begin{exercise}
Why are elementary matrices invertible?
\end{exercise}
\newpage
\begin{exercise}
 Compute $E_1A$, then $E_2E_1A$, then $E_3E_2E_1A$.  What do you notice about the structure of each of the elementary matrices, and the structure of each of the products?
 \end{exercise}
 \vfill
 \begin{exercise}
 Write the resulting upper triangular matrix as the product $L_2L_1A$, where $L_1$ takes $A$ and zeroes out the first column except for the first entry, and $L_2$ takes $L_1A$ and zeros out the second column except for the first two entries.  (These are called elimination matrices.)
 \end{exercise}
\vfill
\newpage
 \begin{exercise}
Do everything again, except instead of using numbers, consider a general matrix 
    \begin{align*}
        A = \begin{bmatrix} a_{11} & a_{12} & a_{13} \\
        a_{21} & a_{22} & a_{23} \\
        a_{31} & a_{32} & a_{33}\end{bmatrix}
    \end{align*}
    and assume $a_{ij} \neq 0$ for all $i,j \in \{1,2,3\}$.
\end{exercise}



\end{document}